\documentclass[UTF8,a4paper,12pt]{ctexart}

\usepackage[left=2.6cm,right=2.6cm,top=2.5cm,bottom=2.5cm]{geometry}
\usepackage{booktabs,tabularx,enumitem,xcolor,amssymb}
\usepackage{hyperref}
\hypersetup{hidelinks}

\setlength{\parindent}{2em}
\setlength{\parskip}{0.3em}
\renewcommand{\arraystretch}{1.4}
\setlist[itemize]{leftmargin=2.2em,itemsep=0.2em,topsep=0.25em}
\setlist[enumerate]{leftmargin=2.4em,itemsep=0.35em,topsep=0.25em}

\newcolumntype{Y}{>{\raggedright\arraybackslash}X}
\newcommand{\Fill}[1]{\textcolor{gray}{\textit{#1}}}
\newcommand{\Field}[1]{\par\noindent\textbf{#1}}

\begin{document}

\begin{center}
  {\Large\bfseries AI 工具使用详情}
\end{center}

\vspace{0.6em}

\noindent\fbox{\parbox{0.965\linewidth}{\small
\textbf{本文件的依据与填写口径。}结构按《全国大学生数学建模竞赛人工智能工具使用规定
（2026 年试行）》第 4 条的四项要求编排，四个小节与四项一一对应；版式参考指导教师建议的
样例。本文件只记录竞赛期间真实发生的使用情况：没用过的工具不列，没发生的交互不编。
文件名须为 \texttt{AI 工具使用详情.pdf}，随支撑材料提交。本文件、论文附录与支撑材料压缩包的
任何位置（含文件夹名、文件名、文档属性）都不得出现队员姓名、学校或赛区信息，也不得写入
账号、API Key 等敏感信息。
}}

\section*{一、AI 工具名称和版本}
\addcontentsline{toc}{section}{一、AI 工具名称和版本}

\begin{center}
  表 1\quad AI 工具名称、版本及使用目的
\end{center}

\noindent\begin{tabularx}{\linewidth}{p{3.2cm}p{3.4cm}Y}
  \toprule
  工具名称 & 版本或型号 & 具体使用目的和环节 \\
  \midrule
  \Fill{实际使用的工具名} & \Fill{版本号、模型名或访问日期} & \Fill{只写本队真实用到的环节} \\
  \Fill{未用到则删除本行} & \Fill{---} & \Fill{---} \\
  \bottomrule
\end{tabularx}

\vspace{0.4em}
\noindent\textit{填写要点：}版本或型号要具体到可复核的粒度（如模型名加日期），不要写“最新版”
或“常见大模型”。同一工具在不同环节使用时合并为一行，用途栏分列环节。

\section*{二、具体使用目的与环节}
\addcontentsline{toc}{section}{二、具体使用目的与环节}

按实际发生的环节逐条写，每条说明“本队先做了什么、把什么交给 AI、AI 承担哪一小段”。
核心建模与分析必须由本队主导，这一点不因某个环节用了 AI 而改变。

\begin{enumerate}
  \item \textbf{\Fill{环节名，如：剖析题目背景}}：\Fill{先独立完成了什么、在哪一步引入 AI、
        给了什么材料、期望 AI 解决什么。}
  \item \textbf{\Fill{环节名，如：代码语法检查与调试}}：\Fill{同上。}
  \item \textbf{\Fill{环节名，如：图表优化}}：\Fill{同上。}
  \item \textbf{\Fill{环节名，如：论文语言润色}}：\Fill{同上。语言润色在第四节可不列采纳记录，
        但本节仍应写明。}
\end{enumerate}

\section*{三、主要提示方式与使用过程说明}
\addcontentsline{toc}{section}{三、主要提示方式与使用过程说明}

\subsection*{主要提示方式}

\begin{itemize}
  \item \Fill{网页对话框提问：输入待解决问题，或复制赛题片段、待润色正文向 AI 提问；}
  \item \Fill{上传文件/图片或粘贴报错日志：把调试中遇到的报错截图或日志发给 AI；}
  \item \Fill{智能体自动执行任务：通过提示词让 AI 直接修改代码或生成图表；}
  \item \Fill{其他实际使用的方式，未用到的删除。}
\end{itemize}

\subsection*{典型交互示例}

每个示例四栏齐全，重点是最后一栏：说明本队采纳、部分采纳还是不采纳，以及依据什么作出该判断。
不采纳的例子同样应当保留，它比“全部采纳”更能说明人工审查真实发生过。

\vspace{0.3em}
\noindent\textbf{示例 1\quad 工具：\Fill{工具名与版本}}
\Field{环节：}\Fill{对应第二节的哪一条}
\Field{提问 / 提示词：}\Fill{原样或摘要照录，不事后美化}
\Field{回答摘要：}\Fill{AI 实际给出的结论要点}
\Field{队伍处理方式：}\Fill{采纳 / 部分采纳 / 不采纳，并写出理由，例如“附件未提供该参数，
故改用可观测量替代，未直接采纳”}

\vspace{0.5em}
\noindent\rule{\linewidth}{0.35pt}

\vspace{0.5em}
\noindent\textbf{示例 2\quad 工具：\Fill{工具名与版本}}
\Field{环节：}\Fill{}
\Field{提问 / 提示词：}\Fill{}
\Field{回答摘要：}\Fill{}
\Field{队伍处理方式：}\Fill{}

\vspace{0.5em}
\noindent\rule{\linewidth}{0.35pt}

\vspace{0.5em}
\noindent\textbf{示例 3\quad 工具：\Fill{工具名与版本}}
\Field{环节：}\Fill{}
\Field{提问 / 提示词：}\Fill{}
\Field{回答摘要：}\Fill{}
\Field{队伍处理方式：}\Fill{}

\section*{四、对 AI 输出的采纳、人工修改与核验的主要情况}
\addcontentsline{toc}{section}{四、对 AI 输出的采纳、人工修改与核验的主要情况}

按工具分组叙述。规定第 4 条第（4）项注明\textbf{语言润色除外}，因此纯润色环节不必在本节
逐条列采纳记录；其余环节都要写清“AI 给了什么、本队改了什么、凭什么确认可用”。

\vspace{0.3em}
\noindent\textbf{工具：\Fill{工具名与版本}}

\Fill{在哪个环节、对哪一处内容，AI 的输出被采纳、被修正还是被否决。修正的要写出改动前后的
差别，例如“AI 建议把窗口收窄到 1600--2400 并由 5 条曲线减为 3 条；观察预览后改为收窄到
1600--2100 并保留 5 条”。被否决的要写出否决依据。}

\Field{核验方式：}\Fill{队内讨论、独立重算、量纲检查、公式推导、与附件数据对账、
代码复跑、图表前后对比、查阅资料等，只写实际做过的。}

\vspace{0.5em}
\noindent\rule{\linewidth}{0.35pt}

\vspace{0.5em}
\noindent\textbf{工具：\Fill{工具名与版本}}

\Fill{同上。}

\Field{核验方式：}\Fill{同上。}

\vspace{0.8em}
\noindent\rule{\linewidth}{0.35pt}

\vspace{0.5em}
\noindent\textbf{总结段（按实际情况改写，不得照抄）}

\Fill{本参赛队使用 AI 工具主要用于【实际用途】。上述 AI 参与完成的内容均已由队员逐项人工
审查与核实，核验方式见各工具条目。【此处必须与事实一致：如果核心模型的建立、推导与求解
确由本队自主完成，写明；如果某一部分由 AI 参与生成，则写明是哪一部分、由谁审查、如何核实，
不要套用“仅作为辅助效率工具”这类与实际不符的表述。】}

\vspace{1.2em}
\noindent\fbox{\parbox{0.965\linewidth}{\small
\textbf{正文侧还须完成两件事，本文件不能替代。}

\textbf{（1）AI 工具使用声明。}按规定第 3 条，在论文\textbf{参考文献之前}设“AI 工具使用声明”，
下列二者\textbf{择一}，不得同时保留：

\begin{itemize}
  \item 未使用：\textbf{本参赛队在竞赛过程中未使用任何 AI 工具。}
  \item 使用：\textbf{本参赛队在竞赛过程中使用了 AI 工具，主要用于【简要用途，如语言润色、
        代码调试等】，详细使用情况见支撑材料。}
\end{itemize}

\textbf{（2）赛区附加要求。}广东赛区《报名和参赛须知》（2026-04-10 修订）另要求：AI 生成的
内容\textbf{在正文相应位置标注}，并\textbf{在参考文献中列出所用 AI 工具}。国规第 6 条写明
“自 2026 年 9 月 1 日起试行，此前相关规定与本规定不一致的以本规定为准”，赛区须知的修订
日期早于该日期，故国规在冲突处优先；但赛区须知同时声明“不遵守本参赛须知的队伍将失去评奖
资格”。两者取并集不违反任何一方，因此\textbf{建议同时做到}：参考文献前放声明、正文标注
AI 生成处、参考文献列出 AI 工具、支撑材料放本 PDF。最终以当届组委会与赛区的书面通知为准。
}}

\end{document}
